% \begin{document}
\chapter{Separation Logic}

So far, we have looked at a version of Hoare logic that supports reasoning about IMP programs in
terms of pre- and post-conditions on stores. In this lecture, we consider a language with pointers
and dynamic memory allocation. We define a big-step operational semantics for the language as a
relation on program states $(\sigma,h)$ with a store and a heap. We show how Hoare logic fails to
support modular reasoning due to aliasing between pointers. Finally, we introduce
\emph{separation logic}, which extends Hoare logic's assertions with predicates on heaps and
provides a \emph{frame} rule that supports modular proofs.

\section{IMP with Heaps}
%
As a first step, we extend IMP with standard commands for manipulating values on the heap. Note: in
this lecture we will write $e$ for arithmetic expressions rather than $a$, as it is more intuitive:
\begin{align*}
  c\ ::=\ & \kw{skip}\mid x:=e\mid c_1;c_2\mid\kw{if}\ b\ \kw{then}\ c_1\ \kw{else}\ c_2\mid\kw{while}\ b\ \kw{do}\ c\\
  \mid\ & x:=\kw{new}(e)\mid\kw{free}(e)\mid x:=\ast e\mid\ast e_1:=e_2
\end{align*}
The new commands in this language include:
\begin{itemize}
  \item $x:=\kw{new}(e)$, which evaluates $e$ to a value $v$, allocates storage for $v$ on the heap,
    and assigns the pointer to that storage to $x$ in the store;
  \item $x:=\ast e$, which evaluates $e$ to a pointer $p$, loads the value $v$ associated with $p$
    from the heap, and assigns $v$ to $x$ in the store;
  \item $\ast e_1:=e_2$, which evaluates $e_1$ to a pointer $p$ and $e_2$ to a value $v$, and then
    stores $v$ at the location associated with $p$ on the heap; and
  \item $\kw{free}(e)$, which evaluates $e$ to a pointer $p$ and deallocates the heap storage
    associated with $p$.
\end{itemize}
%
\section{Big-Step Operational Semantics}
%
To define the semantics of this extension of IMP, we will need states that keep track both of the
local variables as well as the heap. We will model states as pairs $(\sigma,h)$ where
\begin{itemize}
  \item $\sigma\in\mathrm{Var}\to\mathbb{Z}$ is the store, which for simplicity we take to be a
    total function, and
  \item $h\in\mathrm{Addr}\rightharpoonup_{\mathrm{fin}}\mathbb{Z}$ is the heap, a finite partial
    function.
\end{itemize}
Here $\mathrm{Addr}\subseteq\mathbb{Z}\setminus\{0\}$ is the set of addresses. Pointers are just
integers, so they can be stored in variables, stored on the heap, and computed with arithmetic
expressions. We exclude $0$ so that it can serve as the null pointer: no cell is ever stored at
address $0$. As usual, we write $\dom(h)$ for the domain of the heap and $h[p\mapsto v]$ for the
heap that maps $p$ to $v$ and otherwise behaves like $h$.
%
In the formal semantics, commands evaluate in one big step to an outcome $\omega$, which is either
a final state or a memory error:
\[
  \ev{\sigma,h,c}\sem\omega\qquad\qquad \omega\ ::=\ (\sigma',h')\mid\err
\]
while arithmetic and boolean expressions remain pure,
\[
  \ev{\sigma,e}\sem v
\]
%
\paragraph{Basic commands}
\[
  \infr{}{\ev{\sigma,h,\kw{skip}}\sem(\sigma,h)}{Skip}
  \qquad
  \infr{\ev{\sigma,e}\sem v}{\ev{\sigma,h,x:=e}\sem(\sigma[x\mapsto v],h)}{Assign}
\]
\[
  \infr{\ev{\sigma,h,c_1}\sem(\sigma_1,h_1)\qquad\ev{\sigma_1,h_1,c_2}\sem\omega}
       {\ev{\sigma,h,c_1;c_2}\sem\omega}{Seq}
\]
%
\paragraph{Conditionals and loops}
\[
  \infr{\ev{\sigma,b}\sem\mathrm{true}\qquad\ev{\sigma,h,c_1}\sem\omega}
       {\ev{\sigma,h,\kw{if}\ b\ \kw{then}\ c_1\ \kw{else}\ c_2}\sem\omega}{If-True}
  \qquad
  \infr{\ev{\sigma,b}\sem\mathrm{false}\qquad\ev{\sigma,h,c_2}\sem\omega}
       {\ev{\sigma,h,\kw{if}\ b\ \kw{then}\ c_1\ \kw{else}\ c_2}\sem\omega}{If-False}
\]
\[
  \infr{\ev{\sigma,b}\sem\mathrm{false}}
       {\ev{\sigma,h,\kw{while}\ b\ \kw{do}\ c}\sem(\sigma,h)}{While-False}
\]
\[
  \infr{\ev{\sigma,b}\sem\mathrm{true}\qquad\ev{\sigma,h,c}\sem(\sigma_1,h_1)\qquad
        \ev{\sigma_1,h_1,\kw{while}\ b\ \kw{do}\ c}\sem\omega}
       {\ev{\sigma,h,\kw{while}\ b\ \kw{do}\ c}\sem\omega}{While-True}
\]
%
\paragraph{Heap commands}
\[
  \infr{\ev{\sigma,e}\sem p\qquad p\in\dom(h)\qquad h(p)=v}
       {\ev{\sigma,h,x:=\ast e}\sem(\sigma[x\mapsto v],h)}{Load}
  \qquad
  \infr{\ev{\sigma,e_1}\sem p\qquad\ev{\sigma,e_2}\sem v\qquad p\in\dom(h)}
       {\ev{\sigma,h,\ast e_1:=e_2}\sem(\sigma,h[p\mapsto v])}{Store}
\]
\[
  \infr{\ev{\sigma,e}\sem v\qquad p\in\mathrm{Addr}\setminus\dom(h)}
       {\ev{\sigma,h,x:=\kw{new}(e)}\sem(\sigma[x\mapsto p],h[p\mapsto v])}{New}
  \qquad
  \infr{\ev{\sigma,e}\sem p\qquad p\in\dom(h)}
       {\ev{\sigma,h,\kw{free}(e)}\sem(\sigma,h|_{\dom(h)\setminus\{p\}})}{Free}
\]
Here $h|_A$ is the restriction of $h$ to the addresses in $A$. Note that \rn{New} may choose any
unused address $p$, so the semantics is nondeterministic.
%
\paragraph{Memory errors} A load, store, or free at an address that is not allocated produces a
memory error, and errors propagate through sequencing and loops:
\[
  \infr{\ev{\sigma,e}\sem p\qquad p\notin\dom(h)}{\ev{\sigma,h,x:=\ast e}\sem\err}{Load-Err}
  \qquad
  \infr{\ev{\sigma,e_1}\sem p\qquad p\notin\dom(h)}{\ev{\sigma,h,\ast e_1:=e_2}\sem\err}{Store-Err}
\]
\[
  \infr{\ev{\sigma,e}\sem p\qquad p\notin\dom(h)}{\ev{\sigma,h,\kw{free}(e)}\sem\err}{Free-Err}
  \qquad
  \infr{\ev{\sigma,h,c_1}\sem\err}{\ev{\sigma,h,c_1;c_2}\sem\err}{Seq-Err}
\]
\[
  \infr{\ev{\sigma,b}\sem\mathrm{true}\qquad\ev{\sigma,h,c}\sem\err}
       {\ev{\sigma,h,\kw{while}\ b\ \kw{do}\ c}\sem\err}{While-Err}
\]
The rules for \rn{Seq}, \rn{If}, and \rn{While-True} above already pass along an error from their
last premise. Since $0\notin\mathrm{Addr}$, any load, store, or free at address $0$ produces
$\err$. No other command can fail: the store is total, so expressions always evaluate, and we
assume that $\mathrm{Addr}$ is infinite, so \rn{New} can always find an unused address in the
finite heap. Thanks to the outcome $\err$, a program that crashes is now distinguishable from one
that runs forever, which has no derivation at all.
%
\section{The Rule of Constancy}
%
The following rule is admissible in standard Hoare logic:
\[
  \infr{\vdash\triple{P}{c}{Q}\qquad\fv(R)\cap\modv(c)=\emptyset}
       {\vdash\triple{P\wedge R}{c}{Q\wedge R}}{Constancy}
\]
Here $\fv(R)$ are the free variables of $R$ and $\modv(c)$ are the variables that $c$ may modify.
For example, $\fv(y=0)=\{y\}$ and $\modv(x:=x+1)=\{x\}$.
%
Intuitively, the rule of constancy captures a form of local reasoning. It allows us to first prove
a simple Hoare triple and then extend it to a more complicated triple by conjoining the same
predicate to the pre- and post-conditions. For instance, using this rule we can strengthen
\[
  \triple{x>0}{x:=x+1}{x>1}
\]
to
\[
  \triple{x>0\wedge y=0}{x:=x+1}{x>1\wedge y=0}
\]
%
\paragraph{Free and Modified Variables}
The function $\fv(P)$ returns the set of program variables that occur free in assertion $P$.
\begin{align*}
  \fv(n)&=\fv(\mathrm{true})=\fv(\mathrm{false})=\emptyset\\
  \fv(x)&=\{x\}\\
  \fv(a_1+a_2)&=\fv(a_1\times a_2)=\fv(a_1\le a_2)=\fv(a_1)\cup\fv(a_2)\\
  \fv(P\wedge Q)&=\fv(P\vee Q)=\fv(P\Rightarrow Q)=\fv(P)\cup\fv(Q)\\
  \fv(\neg P)&=\fv(\forall i.\,P)=\fv(\exists i.\,P)=\fv(P)
\end{align*}
The quantifier cases do not remove anything, because $i$ is a logical variable, not a program
variable.
%
The function $\modv(c)$ returns the set of program variables modified by command $c$.
\begin{align*}
  \modv(\kw{skip})&=\modv(\ast e_1:=e_2)=\modv(\kw{free}(e))=\emptyset\\
  \modv(x:=e)&=\modv(x:=\ast e)=\modv(x:=\kw{new}(e))=\{x\}\\
  \modv(c_1;c_2)&=\modv(\kw{if}\ b\ \kw{then}\ c_1\ \kw{else}\ c_2)=\modv(c_1)\cup\modv(c_2)\\
  \modv(\kw{while}\ b\ \kw{do}\ c)&=\modv(c)
\end{align*}
%
\paragraph{Issues Related to Aliasing} Unfortunately, the rule of constancy is not sound when we
extend our language of assertions with heap predicates. For example, suppose that we add an
assertion $e_1\pto e_2$, which says that the heap contains a cell at address $e_1$ holding $e_2$
(and possibly other cells as well). Now consider trying to use the rule of constancy to add a
``constant fact'' about the heap. We might first prove the following triple:
\[
  \{\,x\pto-\,\}\ \ast x:=4\ \{\,x\pto 4\,\}.
\]
and then use the rule of constancy to obtain:
\[
  \{\,x\pto-\ \wedge\ y\pto 3\,\}\ \ast x:=4\ \{\,x\pto 4\ \wedge\ y\pto 3\,\}.
\]
Here $\modv(\ast x:=4)=\emptyset$, so the side condition of the rule is satisfied trivially.
However, if $x$ and $y$ are aliases for each other, then the postcondition is false. Of course, it
is possible to complicate the pre- and post-condition to capture disjointness predicates, but this
quickly becomes impractical---in principle, every pointer on the heap might be aliased to every
other pointer, which leads to a combinatorial explosion. Such complications motivated the
development of separation logic.
%
\section{Heap Assertions}
%
Before defining separation logic, let us extend our language of assertions with heap predicates:
\[
  H,J,K\ ::=\ \emp\mid\pure{P}\mid e_1\pto e_2\mid H_1\sep H_2\mid H_1\hand H_2\mid H_1\hor H_2
  \mid\forall i.\,H\mid\exists i.\,H
\]
Intuitively, these predicates can be understood as follows.
\begin{itemize}
  \item $\emp$ : heap is empty,
  \item $\pure{P}$ : heap is empty and $P$ holds,
  \item $e_1\pto e_2$ : heap consists of exactly one cell with pointer $e_1$ and value $e_2$ (this
    is stricter than the informal reading of $\pto$ in the previous section),
  \item $H_1\sep H_2$ : heap splits into two disjoint pieces, one satisfying $H_1$ and the other
    $H_2$,
  \item $H_1\hand H_2$ and $H_1\hor H_2$ are heap assertion versions of the standard boolean
    connectives, and
  \item $\forall i.\,H$ and $\exists i.\,H$ are heap assertion versions of the standard
    quantifiers.
\end{itemize}
More formally, we can model their semantics as follows:
\begin{align*}
  (\sigma,h)&\models_I\emp && \text{if } h=\emptyset\\
  (\sigma,h)&\models_I\pure{P} && \text{if } h=\emptyset\wedge\sigma\models_I P\\
  (\sigma,h)&\models_I e_1\pto e_2 && \text{if } h=\{(p,v)\}\text{ where }\ev{\sigma,e_1}_I\sem p\text{ and }\ev{\sigma,e_2}_I\sem v\\
  (\sigma,h)&\models_I H_1\sep H_2 && \text{if } \exists h_1,h_2.\ h=h_1\uplus h_2\wedge(\sigma,h_1)\models_I H_1\wedge(\sigma,h_2)\models_I H_2\\
  (\sigma,h)&\models_I H_1\hand H_2 && \text{if } (\sigma,h)\models_I H_1\text{ and }(\sigma,h)\models_I H_2\\
  (\sigma,h)&\models_I H_1\hor H_2 && \text{if } (\sigma,h)\models_I H_1\text{ or }(\sigma,h)\models_I H_2\\
  (\sigma,h)&\models_I \forall i.\,H && \text{if } (\sigma,h)\models_{I[i\mapsto v]}H\text{ for all }v\\
  (\sigma,h)&\models_I \exists i.\,H && \text{if } (\sigma,h)\models_{I[i\mapsto v]}H\text{ for some }v
\end{align*}
Here $h_1\uplus h_2$ is the union of heaps $h_1$ and $h_2$ with disjoint domains (it is undefined
if they overlap). Free variables extend to heap assertions in the obvious way, e.g.,
$\fv(e_1\pto e_2)=\fv(e_1)\cup\fv(e_2)$, $\fv(\emp)=\emptyset$, and $\fv(\pure{P})=\fv(P)$.
%
Quantifiers bind logical variables (never program variables), which range over integers, including
addresses, and over finite lists of integers; $I$ maps each to a value of the right kind. We
typically write $i,j,a,b$ for integers and $L$ for lists.
%
We will often abbreviate $e\pto-\ \triangleq\ \exists i.\,e\pto i$.
%
As heaps only contain addresses, $0\pto v$ is unsatisfiable for every $v$: there is no heap whose
single cell is at address $0$. Similarly, predicates like $x\pto 5\sep x\pto 5$ are unsatisfiable
as the separating conjunction operator, $\sep$, enforces disjointness---even when the two cells
would hold the same value.
%
\paragraph{Equivalence and entailment.} Two heap assertions are equivalent, written
$H_1\equiv H_2$, if they hold in exactly the same states:
\[
  H_1\equiv H_2\ \triangleq\ \forall\sigma,h,I.\ (\sigma,h)\models_I H_1\iff(\sigma,h)\models_I H_2.
\]
(Note that it is not enough for $H_1$ and $H_2$ to be both valid or both invalid---for example,
$x\pto 1$ and $x\pto 2$ are both invalid, but they are certainly not equivalent.) Similarly, $H_1$
\emph{entails} $H_2$, written $H_1\entails H_2$, if every state satisfying $H_1$ also satisfies
$H_2$.
%
Separating conjunction satisfies several familiar algebraic laws:
\begin{align*}
  H\sep J&\equiv J\sep H && \text{(commutativity)}\\
  (H\sep J)\sep K&\equiv H\sep(J\sep K) && \text{(associativity)}\\
  H\sep\emp&\equiv H && \text{(unit)}
\end{align*}
However, $\sep$ and $\hand$ behave differently. In general,
$\pure{P}\sep H\not\equiv\pure{P}\hand H$: the left-hand side says that $P$ holds and $H$
describes the whole heap, while the right-hand side requires the heap to be empty. For instance,
$\pure{\mathrm{true}}\sep x\pto 1$ is satisfiable, but $\pure{\mathrm{true}}\hand x\pto 1$ is not.
When $H$ is $\emp$, the two do coincide: $\pure{P}\sep\emp\equiv\pure{P}\hand\emp$. In practice,
$\pure{P}\sep H$ is the idiomatic way to state a pure fact $P$ alongside a heap assertion $H$, and
we will use entailments such as $\pure{P}\sep H\entails H$ freely.
%
\section{Valid Triples}
%
Now we define the notion of valid triples. We use the standard definition for separation logic,
which is partial correctness with \emph{fault avoidance}, following Reynolds and O'Hearn: a valid
triple guarantees that the program never produces a memory error, and that if it terminates, the
postcondition holds. The explicit outcome $\err$ is what makes this possible. If memory errors
simply got stuck, a crashing program would have no derivation, just like a program that runs
forever, and a partial correctness triple would hold vacuously for it. With fault avoidance,
$\triple{\emp}{\kw{free}(x)}{\mathrm{false}}$ is not valid, because $\kw{free}(x)$ produces $\err$
on the empty heap. On the other hand, $\triple{\emp}{\kw{while}\ \mathrm{true}\ \kw{do}\ \kw{skip}}{\mathrm{false}}$
is valid: the loop never crashes and never terminates. As usual for partial correctness,
termination must be proved separately if it is needed.
%
\medskip
\noindent\textbf{Definition (Valid Triple).} A triple is \emph{valid}, written
$\models\triple{H}{c}{J}$, if for all $\sigma,h$, and $I$ with $\sigma,h\models_I H$:
\begin{itemize}
  \item $\ev{\sigma,h,c}\sem\err$ is not derivable, and
  \item every $\sigma',h'$ with $\ev{\sigma,h,c}\sem(\sigma',h')$ satisfies $\sigma',h'\models_I J$.
\end{itemize}
Because \rn{New} makes the semantics nondeterministic, both conditions quantify over all
executions: a valid triple rules out an error for every possible choice of fresh addresses.
%
Note that the precondition $H$ describes the entire heap in which $c$ runs, and $c$ must not
produce an error in any such heap. So a valid triple implicitly says that $c$ only accesses cells
described by $H$---its \emph{footprint}. For example, $\triple{x\pto 3}{\ast x:=4}{x\pto 4}$ is
valid, but $\triple{\emp}{\ast x:=4}{\emp}$ is not, because $\ast x:=4$ produces $\err$ on the
empty heap.
%
\section{The Frame Rule}
%
A major appeal of separation logic is that it supports the so-called frame rule:
\[
  \infr{\vdash\triple{H}{c}{J}\qquad\fv(K)\cap\modv(c)=\emptyset}
       {\vdash\triple{H\sep K}{c}{J\sep K}}{Frame}
\]
The rule captures modular reasoning about the heap. Intuitively, we first prove a specification for
$c$ using only the ``heap'' footprint it needs. Then the command can be framed in any larger
context obtained by combining the footprint with a disjoint heap.
%
\paragraph{Soundness.} Why is the frame rule sound? Write $h\perp h_0$ when $h$ and $h_0$ have
disjoint domains. The proof rests on four facts, the first three about the semantics (each proved
by induction on derivations) and the last about assertions (proved by induction on the assertion):
\begin{itemize}
  \item \emph{Safety monotonicity}: if $\ev{\sigma,h,c}\sem\err$ is not derivable, then neither is
    $\ev{\sigma,h\uplus h_0,c}\sem\err$, for any $h_0\perp h$. (If a load, store, or free finds its
    address in $h$, it also finds it in $h\uplus h_0$.)
  \item \emph{Frame property}: if $\ev{\sigma,h,c}\sem\err$ is not derivable and
    $\ev{\sigma,h\uplus h_0,c}\sem(\sigma',h')$, then $h'=h_1'\uplus h_0$ for some $h_1'$ with
    $\ev{\sigma,h,c}\sem(\sigma',h_1')$. In other words, a run in the larger heap leaves $h_0$
    untouched and behaves like a run in $h$.
  \item \emph{Unmodified variables}: if $\ev{\sigma,h,c}\sem(\sigma',h')$, then
    $\sigma'(y)=\sigma(y)$ for every $y\notin\modv(c)$.
  \item \emph{Coincidence} (also called the agreement or relevance lemma): if $\sigma$ and
    $\sigma'$ agree on $\fv(K)$, then $\sigma,h\models_I K$ if and only if $\sigma',h\models_I K$.
    That is, an assertion depends on the store only through its free variables.
\end{itemize}
The interesting case of the frame property is \rn{New}. The run in $h\uplus h_0$ allocates some
$p\notin\dom(h\uplus h_0)$. Since $\dom(h)\subseteq\dom(h\uplus h_0)$, this $p$ is also unused in
$h$, so the run in $h$ alone can pick the same $p$, and the results match:
$(h\uplus h_0)[p\mapsto v]=h[p\mapsto v]\uplus h_0$. This is exactly why \rn{New} is allowed to
choose any unused address. With a deterministic allocator, say $p=1+\max\dom(h)$, the two runs
could allocate different addresses, and the frame rule would be unsound: on the empty heap the
allocator returns $1$, so $\triple{\emp}{x:=\kw{new}(1)}{\pure{x=1}\sep x\pto 1}$ would be valid,
but framing in $1\pto 5$ gives a false triple, because the allocator would now return $2$.
%
\medskip
\noindent\textbf{Theorem (Soundness of Frame).} If $\models\triple{H}{c}{J}$ and
$\fv(K)\cap\modv(c)=\emptyset$, then $\models\triple{H\sep K}{c}{J\sep K}$.
%
\medskip
\noindent\emph{Proof.} Suppose $\sigma,h\models_I H\sep K$. Then $h=h_1\uplus h_0$ with
$\sigma,h_1\models_I H$ and $\sigma,h_0\models_I K$. By the premise, $\ev{\sigma,h_1,c}\sem\err$ is
not derivable, so by safety monotonicity neither is $\ev{\sigma,h,c}\sem\err$. Now suppose
$\ev{\sigma,h,c}\sem(\sigma',h')$. By the frame property, $h'=h_1'\uplus h_0$ with
$\ev{\sigma,h_1,c}\sem(\sigma',h_1')$, so $\sigma',h_1'\models_I J$ by the premise. By the
unmodified-variables fact, $\sigma'$ agrees with $\sigma$ on $\fv(K)$, since
$\fv(K)\cap\modv(c)=\emptyset$; so by coincidence we still have $\sigma',h_0\models_I K$. Hence
$\sigma',h'\models_I J\sep K$. \hfill$\blacksquare$
%
\paragraph{Why fault avoidance?} The proof uses the absence of errors in an essential way. Suppose
instead that memory errors simply got stuck, and triples were ordinary partial correctness
statements. Then $\triple{\emp}{\kw{free}(x)}{\emp}$ would be valid vacuously, since
$\kw{free}(x)$ has no result on the empty heap. But framing in $x\pto 3$ would give
$\triple{x\pto 3}{\kw{free}(x)}{x\pto 3}$, which is false: the free succeeds and the cell is gone.
Fault avoidance rules out the first triple, because it forces a command to own every cell it
touches.
%
\paragraph{Baking in the frame rule.} Some modern presentations, such as Chargu\'eraud's, take a
different approach and build the frame rule into the definition of triples. There,
$\triple{H}{c}{J}$ is defined to be valid if, for every $K$ with $\fv(K)\cap\modv(c)=\emptyset$,
the triple $\triple{H\sep K}{c}{J\sep K}$ is valid in the sense above. The frame rule is then sound
by construction, at the cost of a less direct definition.
%
\section{Small Axioms for Heap Commands}
%
We can define ``small'' axioms for the heap-manipulating commands as follows.
\[
  \infr{x\notin\fv(e_1)\cup\fv(e_2)}
       {\{\,e_1\pto e_2\,\}\ x:=\ast e_1\ \{\,\pure{x=e_2}\sep e_1\pto e_2\,\}}{Load}
  \qquad
  \infr{}{\{\,e_1\pto-\,\}\ \ast e_1:=e_2\ \{\,e_1\pto e_2\,\}}{Store}
\]
\[
  \infr{x\notin\fv(e)}{\{\,\emp\,\}\ x:=\kw{new}(e)\ \{\,x\pto e\,\}}{Alloc}
  \qquad
  \infr{}{\{\,e\pto-\,\}\ \kw{free}(e)\ \{\,\emp\,\}}{Free}
\]
In \rn{Alloc}, the program variable $x$ ends up holding the address of the new cell, so
$x\pto e$ describes that cell. The side condition is needed because $e$ is interpreted in the final
store, where $x$ has been overwritten.
%
A full treatment of separation logic has other proof rules similar to Hoare logic. In the examples
below we also use the rules for sequencing and consequence, where consequence is stated using
entailment ($\entails$) between heap assertions, and the familiar invariant rule for loops:
\[
  \infr{\{\pure{b}\sep I\}\ c\ \{I\}}
       {\{I\}\ \kw{while}\ b\ \kw{do}\ c\ \{\pure{\neg b}\sep I\}}{While}
\]
Note how the loop test is combined with the invariant using $\sep$: $\pure{b}\sep I$ says that $b$
holds and $I$ describes the heap. Using $\pure{b}\hand I$ instead would be wrong: it would also
require an empty heap.
%
\section{Example: Swap}
%
The following program swaps the contents of the cells at $p$ and $q$:
\[
  t:=\ast p;\ u:=\ast q;\ \ast p:=u;\ \ast q:=t
\]
A natural specification is the following, where $a$ and $b$ are logical variables:
\[
  \{\,p\pto a\sep q\pto b\,\}\ \mathrm{swap}\ \{\,p\pto b\sep q\pto a\,\}.
\]
We can prove it using the small axioms, framing away the cells (and pure facts) each command does
not touch. Each use of \rn{Frame} is allowed, since the frame never mentions the assigned variable:
\begin{align*}
  &\{\,p\pto a\sep q\pto b\,\}\\
  &t:=\ast p;\qquad \text{\rn{Load}, framing } q\pto b\\
  &\{\,\pure{t=a}\sep p\pto a\sep q\pto b\,\}\\
  &u:=\ast q;\qquad \text{\rn{Load}, framing } \pure{t=a}\sep p\pto a\\
  &\{\,\pure{t=a}\sep\pure{u=b}\sep p\pto a\sep q\pto b\,\}\\
  &\ast p:=u;\qquad \text{\rn{Store}, framing everything but } p\pto a\\
  &\{\,\pure{t=a}\sep\pure{u=b}\sep p\pto u\sep q\pto b\,\}\\
  &\ast q:=t\qquad \text{\rn{Store}, framing everything but } q\pto b\\
  &\{\,\pure{t=a}\sep\pure{u=b}\sep p\pto u\sep q\pto t\,\}\ \entails\\
  &\{\,p\pto b\sep q\pto a\,\}
\end{align*}
(Before each \rn{Store} we also weaken $p\pto a$ to $p\pto-$, and similarly for $q$.)
%
\paragraph{Is this the only specification?} No! Because $\sep$ forces the two cells to be
disjoint, the specification above says nothing about the case where $p$ and $q$ are aliases. In
that case the program is the identity, which we can specify as follows:
\[
  \{\,\pure{p=q}\sep p\pto a\,\}\ \mathrm{swap}\ \{\,\pure{p=q}\sep p\pto a\,\}.
\]
%
\section{Linked Lists}
%
\paragraph{Representation.} We represent a linked list as a sequence of nodes. The node at address
$p$ occupies two consecutive cells: its value is stored at $p$ and the address of the next node is
stored at $p+1$. We use $0$ as the null pointer; recall that $0\notin\mathrm{Addr}$, so $0$ is
never the address of a cell. For example, the following heap holds the list $[1,2,3]$ at address
$6$:
%
\begin{center}
\begin{tikzpicture}[>=Stealth,
    cell/.style={draw,minimum width=9mm,minimum height=7mm,inner sep=0pt},
    addr/.style={font=\footnotesize}]
  \node[addr] at (-1.5,0) {$p=6$};
  \draw[->] (-1.0,0) -- (-0.1,0);
  % node at 6
  \node[cell] (a1) at (0.2,0) {1};
  \node[cell] (a2) at (1.1,0) {30};
  \node[addr,above=0pt of a1] {6};
  \node[addr,above=0pt of a2] {7};
  % node at 30
  \node[cell] (b1) at (3.2,0) {2};
  \node[cell] (b2) at (4.1,0) {14};
  \node[addr,above=0pt of b1] {30};
  \node[addr,above=0pt of b2] {31};
  % node at 14
  \node[cell] (c1) at (6.2,0) {3};
  \node[cell] (c2) at (7.1,0) {0};
  \node[addr,above=0pt of c1] {14};
  \node[addr,above=0pt of c2] {15};
  \draw[->] (a2.south) -- ++(0,-0.6) -| (b1.south);
  \draw[->] (b2.south) -- ++(0,-0.6) -| (c1.south);
\end{tikzpicture}
\end{center}
%
As a finite map, it is $h=\{6\mapsto 1,\ 7\mapsto 30,\ 30\mapsto 2,\ 31\mapsto 14,\ 14\mapsto 3,\ 15\mapsto 0\}$.
Nodes need not be adjacent or in order; only the pointers matter. (Our $\kw{new}(e)$ allocates a
single cell. A real language would allocate a whole node at once, but we will not need allocation
below.)
%
\paragraph{A list predicate.} To reason about lists, we define a heap predicate
$\mathrm{list}(p,L)$, meaning ``$p$ points to a list holding the mathematical list $L$,'' by
recursion on $L$:
\begin{align*}
  \mathrm{list}(p,[\,])&\ \triangleq\ \pure{p=0}\\
  \mathrm{list}(p,v\ld L)&\ \triangleq\ \exists q.\ p\pto v\sep(p+1)\pto q\sep\mathrm{list}(q,L)
\end{align*}
Unfolding the definition step by step:
\begin{align*}
  &\mathrm{list}(6,[1,2,3])\\
  \equiv\ &\exists q_1.\ 6\pto 1\sep 7\pto q_1\sep\mathrm{list}(q_1,[2,3])\\
  \equiv\ &\exists q_1,q_2.\ 6\pto 1\sep 7\pto q_1\sep q_1\pto 2\sep(q_1+1)\pto q_2\sep\mathrm{list}(q_2,[3])\\
  \equiv\ &\exists q_1,q_2,q_3.\ 6\pto 1\sep 7\pto q_1\sep q_1\pto 2\sep(q_1+1)\pto q_2\sep q_2\pto 3\sep(q_2+1)\pto q_3\sep\pure{q_3=0}
\end{align*}
The heap above satisfies the last line, with $q_1=30$, $q_2=14$, and $q_3=0$.
%
\paragraph{What $\sep$ rules out.} Because each unfolding claims new cells that are disjoint from
all the others, and the list must end at $0$, no list $L$ makes $\mathrm{list}(p,L)$ describe a
cyclic structure. Similarly, $\mathrm{list}(p_1,L_1)\sep\mathrm{list}(p_2,L_2)$ guarantees that the
two lists share no nodes: a shared tail would have to belong to both halves of the heap. In Hoare
logic we would need explicit ``acyclic'' and ``disjoint'' predicates to say these things. Here
$\sep$ gives us both for free.
%
\section*{Further Reading}
%
Interested students are referred to two canonical treatments: Reynolds's lecture notes, or
Chargu\'eraud's habilitation thesis or textbook \emph{Separation Logic Foundations} (Software
Foundations, Volume 6), which this lecture is loosely based on.
%
% \end{document}
